\clearpage
\section{GeTe and SnSe: PBE-Based Selection}
\label{sec:gete-snse-pbe-selection}

Monolayer GeTe and SnSe extend the present 2D GW work to a new pair of
applications.  This subsection preserves the initial 6 September
pseudopotential (PP) and numerical atomic orbital (NAO) screening, before
the subsequent EXX failure and library-basis control described below.
Figure~\ref{fig:gete-snse-pbe-six-panel}
preserves that six-panel PBE comparison against FHI-aims all-electron
(AE) bands, without spin--orbit coupling (SOC).

\paragraph{Comparison and observed gaps.}
Each material uses the same geometry and plotting path across methods, with
each spectrum independently aligned to its path valence-band maximum (VBM).
The PP self-consistent meshes are \(12\times12\times1\); the AE meshes are
\(24\times24\times1\) for GeTe and \(20\times20\times1\) for SnSe.
Thus this is not an identical-mesh comparison.  The AE numerical controls
passed their recorded PBE checks, but remain provisional references, not
finally converged AE benchmarks.  Define the path gap as
\(E_g=\min_{\bk\in\mathrm{path}}E_c(\bk)
-\max_{\bk\in\mathrm{path}}E_v(\bk)\), using the lowest empty and highest
occupied bands, and \(\Delta E_g=E_g^{\rm NAO}-E_g^{\rm AE}\).

\begin{table}[!ht]
  \centering
  \caption{No-SOC PBE path gaps.  These are not global Brillouin-zone
  extrema or absolute band offsets.  Extra printed digits identify the
  frozen data, not a convergence uncertainty.}
  \label{tab:gete-snse-pbe-gaps}
  \small
  \begin{tabular}{llrrr}
    \toprule
    Material & PP/NAO & AE (eV) & NAO (eV) & \(\Delta E_g\) (meV) \\
    \midrule
    GeTe & GTH            & 0.868309 & 1.030718 & +162.409 \\
    GeTe & SG15 v1.0      & 0.868309 & 0.871440 &   +3.131 \\
    GeTe & Dojo stringent & 0.868309 & 0.873450 &   +5.141 \\
    SnSe & GTH            & 0.966363 & 0.903439 &  -62.924 \\
    SnSe & SG15 v1.0      & 0.966363 & 0.955169 &  -11.194 \\
    SnSe & Dojo stringent & 0.966363 & 0.960081 &   -6.282 \\
    \bottomrule
  \end{tabular}
\end{table}
\FloatBarrier

\paragraph{Initial selection (6 September 2026; historical).}
Dojo stringent with the matched TZDP basis was initially selected for both
materials, and GeTe/GTH was rejected for this application.
The small Dojo path-gap differences and the visible near-edge band agreement
motivate this practical PBE-based choice.  The selected orbital sets have
\(6s6p3d2f\) channels and 10-bohr cutoff radii for each element, with 212
NAOs per four-atom cell.  This is a user-directed selection, not a claim of
uniform empty-band accuracy or complete PP/basis validation.  In particular,
the SnSe/SG15 v1.0 panel contains extra low-lying empty branches despite a
close gap; their PP-versus-basis origin is unresolved, and no general
rejection of the SG15 family is inferred.

\paragraph{Waived screening and retained GW requirements.}
Further alternative-PP screening, same-PP plane-wave (own-PW) checks, and
PP/AE head-response selection tests are stopped by user direction:
\textbf{waived, not passed}.  No own-PW or head-selection test is supplied
by this comparison.  This waiver neither establishes that an AE response
calculation is impossible nor disables the production correction:
\textbf{the strict-2D GW \mbox{head/wing} treatment remains enabled}.
Its inputs and numerical validity, the matched auxiliary basis, and GW
convergence still require validation.  This initial PBE comparison alone
does not establish the validity of either material's GW spectrum.

\clearpage
\begin{figure}[!ht]
  \centering
  \includegraphics[width=\textwidth]{figures/gete_snse_pbe_selection_20260906/selected_footer_v2/ae_nao_pbe.pdf}
  \caption[GeTe/SnSe six-panel no-SOC PBE comparison]{No-SOC PBE bands of
  GeTe (top) and SnSe (bottom), comparing GTH, SG15 v1.0, and Dojo stringent
  (left to right).  FHI-aims AE/ABACUS NAO curves are dashed/solid;
  occupied/empty states are red/cyan, with independent path-VBM alignment
  and native band order.  Dojo stringent TZDP was the initial selection on
  this PBE evidence; the later library-basis control is described separately.
  AE references remain provisional;
  no GW result is shown.}
  \label{fig:gete-snse-pbe-six-panel}
\end{figure}
\FloatBarrier

\clearpage
\subsection{GeTe: Library-Basis GW Result}
\label{sec:gete-library-gw}

The initial large-basis route gave abnormal EXX bands even when its final
auxiliary Coulomb matrices were positive.  A positive Coulomb metric alone
does not validate the local RI fit or its contraction with the density
matrix.  The following control instead uses the vetted Dojo-NC-SR library
PP/NAO pair.  Because the PP, physical NAOs and auxiliary basis all changed
relative to the initial screening route, the improvement cannot be
attributed exclusively to an orbital-generation parameter or inverse cutoff.

\paragraph{Calculation and scope.}
The system is a four-atom, puckered orthorhombic GeTe monolayer, without SOC,
at the literature geometry used in the PBE comparison.  The SCF/GW mesh is
\(12\times12\times1\), the cell height is 20~\AA, and the plotted path is
\(\Gamma\)--X--S--Y--\(\Gamma\), with 50 intervals per segment.
The physical TZDP basis has \(3s3p3d2f\) radial multiplicities, a 10-bohr
cutoff and 164 orbitals per cell.  One additional \(g\) function per atom
is used only to generate the auxiliary basis, not to change the physical
one-electron Hilbert space.  Auxiliary generation uses PCA \(10^{-4}\);
the 1508 raw auxiliary functions are compressed to 1050 active functions.
No positive Cs inverse-eigenvalue threshold is applied.
Full strict-2D Coulomb matrices use the direct mixed-Fourier construction
(110~Ry, 64 out-of-plane quadrature points, Gamma order 8), with shrink and
symmetry enabled.  The GW calculation uses 16 minimax frequencies,
12-parameter Pad\'{e} continuation and analytic Gamma head/wing treatment.
These are the tested operating parameters, not a completed GeTe convergence
study.

\begin{table}[!ht]
  \centering
  \caption{GeTe no-SOC path gaps for the library-basis control.  Values are
  computed from the raw occupied and empty manifolds, before display
  reordering.  Edge coordinates are fractional reciprocal coordinates
  \((k_x,k_y)\) in the present crystal axes, with \(k_z=0\).
  EXX is the diagonal EXX@PBE spectrum, not a self-consistent
  Hartree--Fock calculation.}
  \label{tab:gete-library-gw-gaps}
  \begin{tabular}{lccr}
    \toprule
    Method & VBM & CBM & Path gap (eV) \\
    \midrule
    PBE & (0, 0.38) & (0.41, 0) & 0.86874 \\
    EXX@PBE & --- & --- & 3.90577 \\
    \(G_0W_0\)@PBE & (0, 0.39) & (0, 0.40) & 1.38070 \\
    \(G_0W_0\) minus PBE & --- & --- & +0.51196 \\
    \bottomrule
  \end{tabular}
\end{table}
\FloatBarrier

\paragraph{Numerical checks and interpretation.}
All 49 direct full-Coulomb matrices are finite, Hermitian and positive
within the recorded numerical checks; their global minimum eigenvalue is
\(2.51\times10^{-5}\)~Ha.  The full \(201\times164\) GW table is finite,
with no high-empty-state intrusion into the occupied manifold.  The
imaginary-axis self-energy is finite, and the fixed PBE frontier indices
30/31 reproduce the occupied/empty envelope gap.  Thus this calculation
does not reproduce the catastrophic EXX failure of the original route.
It gives a positive GW opening of 0.512~eV and a near-direct path gap: its
sampled VBM and CBM lie at adjacent points along Y--\(\Gamma\), not exactly
the same k point.  This does not establish an exactly direct or global
Brillouin-zone gap.  A GeTe-specific mesh/basis/frequency convergence study
and a same-chain full-q symmetry control remain absent.

\paragraph{Comparison with published work.}
The present geometry comes from the puckered, not hexagonal, monolayer of
Sinambela \emph{et al.}, whose GGA/HSE gaps are 0.90/1.30~eV;
the HSE value is not a GW benchmark~\cite{sinambela2021gete}.
Wang \emph{et al.} report a 0.87~eV PBE indirect gap for the same
orthorhombic phase and explicitly compare GW with and without SOC in
their supplementary Fig.~S1~\cite{wang2021tellurides}.
These results support the present PBE scale.  For the fundamental
quasiparticle gap, Xu \emph{et al.} report \(G_0W_0\)@PBE
\textbf{1.21~eV} in Table~I~\cite{xu2017monochalcogenides}, versus the
present \textbf{1.38070~eV}, a difference of \textbf{+0.17070~eV}.
Their phase and starting functional match, but the relaxed structure,
integration mesh and basis differ, and their methods section does not
explicitly settle SOC or the strict-2D Coulomb boundary.  Their
direction-resolved direct transitions and BSE optical gaps must not be
substituted for the fundamental QP gap.

There is also a \textbf{valley-order difference}: after accounting for
the exchanged crystal-axis labels, the published GW VBM and CBM lie on
perpendicular directions, whereas the present GW extrema both lie along
Y--\(\Gamma\).  This difference is already present in the raw energies
and is not caused or cured by display sorting.  Defining
\(\Delta_c=\min_{Y\Gamma}E_c-\min_{\Gamma X}E_c\), the raw data give
\(+48.08\)~meV at PBE and \(-32.84\)~meV at GW.  The GW valley ordering
therefore reverses: relative to each method's VBM, the
\(\Gamma\)--X/Y--\(\Gamma\) conduction minima are
0.86874/0.91682~eV at PBE and 1.41354/1.38070~eV at GW.
The relative valley separation changes by \(-80.92\)~meV, including the
small shifts of the sampled minima; it is not a rigid gap opening.
The competing GW \(\Gamma\)--X minimum is at (0.42, 0), only
32.84~meV above the Y--\(\Gamma\) CBM at (0, 0.40).
The ordering thus rests on a small energy difference whose convergence is not yet
established.  The present comparison supports the energy scale and a
positive GW gap opening, but does \emph{not} establish quantitative
agreement or a validated valley order.  The competing conduction valleys
and their numerical stability are the next physics question, rather than
further cosmetic changes to the plot.
The two cited theoretical studies are not a proof of universal agreement:
the older values quoted again by Wang \emph{et al.} are not an additional
independent calculation.  An error in the present calculation remains
possible, but the unmatched numerical and structural definitions do not
identify its cause.  We retain this result and its disagreement as
provisional evidence, not as a validated application benchmark.

\begin{figure}[!ht]
  \centering
  \includegraphics[width=0.85\textwidth]{figures/gete_gw_20260909/GeTe_PBE_GW_reordered.pdf}
  \caption[GeTe library-basis PBE/GW bands]{GeTe PBE (cyan dashed) and
  \(G_0W_0\)@PBE (vermilion solid) bands, each aligned to its own occupied
  path VBM.  Cyan hollow circles locate the PBE VBM/CBM and vermilion
  hollow diamonds locate the GW VBM/CBM at their actual sampled k points;
  leaders move only the labels.  The two VBM markers nearly overlap.
  At each original k point, the occupied and empty energies
  are ranked separately within the fixed displayed raw-band windows
  (PBE 25--37; GW 25--36, with occupied bands ending at 30).
  No energies are changed, no interpolation or smoothing is applied, and
  the gaps in Table~\ref{tab:gete-library-gw-gaps} are invariant.  This is
  an energy-ranked dispersion display, not wavefunction-character tracking;
  degeneracies and genuine slope changes are not removed.  Raw spectra,
  the previous raw-index plot and per-k raw-index mappings are retained.
  The result passes numerical sanity checks but is not yet a converged
  quantitative benchmark.}
  \label{fig:gete-library-gw}
\end{figure}
\FloatBarrier
